\section{From Short Tails to Bounded Depth}
\label{sec:depth}

A hole with a short rejected tail can still be arbitrarily long. The second bound
measures its complexity by the depth of a canonical hierarchy of block decompositions, and
shows that this depth is bounded in terms of the tail. The hierarchy is the device that
makes comparisons tractable: each of its levels collects the starting points of the
comparisons that end at one particular cut.

\subsection{The Canonical Desubstitution Hierarchy}

Let $\Gamma$ be an alphabet and $u \in \Gamma^+$ a nonempty word with last letter $\ell$,
the \emph{marker}. Then $u$ factors uniquely as
\[
  u = X_0\,\ell\,X_1\,\ell \cdots X_{h-1}\,\ell \qquad\text{with } X_i \in (\Gamma\setminus\{\ell\})^*.
\]
The \emph{desubstitution} of $u$ is the word $[X_0][X_1]\cdots[X_{h-1}]$ of length $h$
over the alphabet of \emph{blocks} $(\Gamma\setminus\{\ell\})^*$: each block $[X_i]$, a word
over $\Gamma$, is a single letter of the new word. Expanding every block $[X]$ to $X\ell$
recovers~$u$. The
name follows the use of desubstitution and derived sequences in symbolic dynamics
\citep{Durand98}.

\begin{definition}[Hierarchy and depth]\label{def:hierarchy}
The \emph{canonical hierarchy} of a nonempty word $w$ is the sequence $u_0 = w$ and
$u_{j+1} = $ the desubstitution of $u_j$, with marker $\ell_j$, the last letter of $u_j$. A
word is \emph{unary} if all its letters are equal; this notion concerns words and is
unrelated to the arity of $f$. The \emph{depth} of $w$ is the least $j$
such that $u_j$ is unary.
\end{definition}

If $u_j$ is not unary, then $\len{u_{j+1}} < \len{u_j}$, because some letter of $u_j$ differs
from the marker. Hence the depth is well defined and $\depth(w) \le \len{w} - 1$.

\begin{example}[A hierarchy of depth three]\label{ex:hierarchy}
Let $w = 01001101$. Its marker is $1$, and $w = (0)\,1\,(00)\,1\,(\eps)\,1\,(0)\,1$, so
$u_1 = \alpha_1\,\beta_1\,\gamma_1\,\alpha_1$ with the letters $\alpha_1 = [0]$, $\beta_1 = [00]$, and $\gamma_1 = [\eps]$. The marker of
$u_1$ is $\alpha_1$, and $u_1 = (\eps)\,\alpha_1\,(\beta_1\gamma_1)\,\alpha_1$, so $u_2 = \alpha_2\,\beta_2$ with $\alpha_2 = [\eps]$ and $\beta_2 = [\beta_1\gamma_1]$.
Finally, $u_2 = (\alpha_2)\,\beta_2$ gives $u_3 = \alpha_3$ with $\alpha_3 = [\alpha_2]$, which is unary. Thus
$\depth(w) = 3$, with lengths $8,4,2,1$.
\end{example}

The hierarchy is compatible with comparisons: a comparison that does not end at the last
cut of a word is a comparison of the next level, refined by a common suffix of two blocks.

\begin{lemma}[Alignment]\label{lem:alignment}
Let $u$ be the expansion of a nonempty block word $v = [X_0]\cdots[X_{h-1}]$ with marker
$\ell$, and let $s < e < \len{u}$. Then $(s,e)$ is a comparison of $u$ iff $s$ and $e$ lie in
the expansions $X_{i_s}\ell$ and $X_{i_e}\ell$ of two blocks with $i_s < i_e$, at the same
distance from the markers that end these blocks, the block suffixes beginning at these cuts
are equal, and $(i_s+1, i_e+1)$ is a comparison of~$v$.
\end{lemma}

\begin{proof}
Let $e'$ be the first marker position at or after $e$; it exists because $u$ ends with
$\ell$. If $(s,e)$ is a comparison, then $u[s+i] = u[e+i]$ for $0 \le i \le e'-e$, and since no
block contains $\ell$, $s' = s + (e'-e)$ is the first marker position at or after $s$. So $s$
and $e$ lie at the same distance from the markers ending their blocks, the letters between
agree, and $i_s < i_e$ because $s' < e'$. The remaining suffixes $\wsuf{u}{s'+1}$ and
$\wsuf{u}{e'+1}$ are the expansions of $[X_{i_s+1}]\cdots$ and $[X_{i_e+1}]\cdots$, and the
second is a prefix of the first. An expansion determines its block word, since the markers
delimit the blocks, so the block word after $i_e$ is a prefix of the block word after $i_s$:
$(i_s+1,i_e+1)$ is a comparison of $v$. The converse reverses these steps.
\end{proof}

\subsection{Cores and Supports}

We now relate the levels of the hierarchy to the cuts of the original word. Every position
of $u_j$ carries a \emph{core}, which is either an original cut of $w$ (an \emph{ordinary}
core) or a symbolic \emph{start core}. At level $0$, position $i$ has core $i$. At level
$j+1$, the first position has a new start core $\star_{j+1}$, and position $i \ge 1$ inherits
the core of the $(i-1)$-st occurrence of the marker $\ell_j$ in $u_j$, counting occurrences
from~$0$. For a level $j$, let $\Cores_j$ be the set of ordinary cores of $u_j$.

The core of a block is thus the core of the marker that precedes it, and the block
is read from that marker. A block $[X]$ therefore has two expansions. Its
\emph{ordinary expansion} $X\ell$ places the marker after the block and
reconstructs the word. Its \emph{shifted expansion} $\ell X$ places the preceding
marker before the block: it is the factor of $w$ that starts at the core of the
block, read from the reader state at that core. The first block has no preceding
marker, which is why it receives a symbolic start core. Iterating the shifted
expansion down to level zero gives a map $\xi_j$ from level-$j$ letters to words:
$\xi_0(a) = a$, and $\xi_{j+1}([X])$ is $\xi_j(\ell_j)$ followed by the images
under $\xi_j$ of the letters of~$X$.

\begin{lemma}[Cores]\label{lem:cores}
Let $w$ be nonempty.
\begin{enumerate}[leftmargin=*]
\item $\Cores_{j+1} \subseteq \Cores_j$. The start cores of $u_j$ occupy an initial segment
  of its positions, are pairwise distinct, and number at most $j$; the ordinary cores
  increase from left to right.
\item If $\Cores_j \ne \emptyset$ and $t_j = \max \Cores_j$, then
  \[
    \Cores_{j+1} = \{\, s \mid s < t_j \text{ and } \wsuf{w}{t_j} \pre \wsuf{w}{s} \,\}.
  \]
  In particular, $(s,t_j)$ is a comparison of $w$ for every $s \in \Cores_{j+1}$. If
  $\Cores_j = \emptyset$, then $\Cores_{j+1} = \emptyset$.
\end{enumerate}
\end{lemma}

The cut $t_j$ is the core of the last position of $u_j$, the \emph{terminal core}, and we
call $\wsuf{w}{t_j}$ the \emph{terminal core word} of level $j$. Part~2 says that the next
level retains exactly the earlier occurrences of the terminal core word. Hence each level of the hierarchy
collects the left endpoints of the comparisons whose right endpoint is the terminal core of
the previous level.

The first two levels of \cref{ex:hierarchy} are aligned below. The upper rows
list the cuts of $w$ and the letter after each cut; the markers $\ell_0=1$ are in
bold. The lower rows place each letter of $u_1$ under its shifted expansion,
which starts at its core.
\begin{center}
\small
\begin{tabular}{@{}l*{9}{c}@{}}
\toprule
cut of $w$ & $0$ & $1$ & $2$ & $3$ & $4$ & $5$ & $6$ & $7$ & $8$\\
letter after the cut & $0$ & $\mathbf 1$ & $0$ & $0$ & $\mathbf 1$ & $\mathbf 1$ & $0$ & $\mathbf 1$ & \\
\midrule
letter of $u_1$ & $\alpha_1$ & \multicolumn{3}{c}{$\beta_1$} & $\gamma_1$
  & \multicolumn{2}{c}{$\alpha_1$} & last $\ell_0$ & \\
\cmidrule(lr){2-2}\cmidrule(lr){3-5}\cmidrule(lr){6-6}\cmidrule(lr){7-8}\cmidrule(lr){9-9}
core & $\star_1$ & $1$ & & & $4$ & $5$ & & $t_0=7$ & \\
\bottomrule
\end{tabular}
\end{center}
The markers at the cuts $1,4,5$ are the occurrences $0,1,2$ of $\ell_0$, so they
become the cores of the positions $1,2,3$ of $u_1$; the first position receives
$\star_1$. Thus $\Cores_0=\{0,\dots,7\}$ and $\Cores_1=\{1,4,5\}$. The shifted
expansions are $\xi_1([X])=1X$, and indeed $\wsub{w}{1}{4}=100=\xi_1(\beta_1)$
and $\wsub{w}{4}{5}=1=\xi_1(\gamma_1)$, whereas the ordinary expansion of
$\beta_1=[00]$ is $001$. The first letter has no preceding marker and covers only
$\wsub{w}{0}{1}=0$. The last marker is not inherited: its cut is the terminal core
$t_0=7$. The terminal core is the original cut attached to the last letter of a
level; at level zero it is $7$, whereas the terminal cut of $w$ is $8$. The terminal
core word $\wsuf{w}{7}=1$ occurs exactly at the earlier cuts $1,4,5$, as Part~2
predicts. At level $1$, the terminal core is $t_1=5$, and
$\wsuf{w}{5}=101=\xi_1(\alpha_1)\,\wsuf{w}{t_0}$ occurs nowhere earlier. Accordingly,
$u_2=\alpha_2\beta_2$ has cores $\star_2,\star_1$: its second position inherits
$\star_1$ from the first occurrence of $\alpha_1$, so $\Cores_2=\emptyset$.

The proof of \cref{lem:cores} is by induction on $j$ and is given in \cref{app:cores}. Its
main invariant states that $\xi_j$ is injective and that the factor of $w$ from an
ordinary core of level $j$ to the next core is the image under $\xi_j$ of the letter at
that core. For the last core, set $t_{-1}=|w|$, so that $w[t_{-1}:]=\eps$; then
$\wsuf{w}{t_j} = \xi_j(\ell_j)\,\wsuf{w}{t_{j-1}}$. The injectivity of $\xi_j$
together with \cref{lem:alignment} identifies the earlier occurrences of the terminal core
word with the cores of the next level.

\subsection{The Return Gap}
Let $(\mathcal M,\Adm,\Rej)$ be any $N$-state reader. For a rejected
word $w$, let $J$ be its first rejected cut and $r=|w|-J$.
Its level-$j$ \emph{state support} is $\Supp_j=\{q_i:i\in\Cores_j\}$, the set of
original reader states at the ordinary cores. These supports decrease by \cref{lem:cores}.

\begin{lemma}[Return gap]\label{lem:return-gap}
If $(s,e)$ is a comparison of a rejected word $w$ and $q_s=q_e$, then
$e-s\le r$.
\end{lemma}
\begin{proof}
Put $m=|w|$ and $\varpi=e-s$. The comparison implies
$w[:s]w[e:]=w[:m-\varpi]$. Deleting the equal-state loop therefore gives
$q_{m-\varpi}=q_m\in\Rej$. By minimality of $J$, $m-\varpi\ge J$,
so $\varpi\le m-J=r$.
\end{proof}
This holds for every reader, without the rejected-path property and without
assuming that $w$ is a hole. In particular, tail zero excludes equal-state comparisons.
\Cref{fig:return-gap} shows the loop deletion and the resulting rejected prefix.
\begin{figure}[t]
\centering
\begin{tikzpicture}[x=9mm,y=7mm,>=Stealth,font=\small]
  \node[anchor=east] at (-0.3,0) {$w$};
  \fill[black!12] (3,-0.16) rectangle (6,0.16);
  \draw[thick] (0,0) -- (12,0);
  \foreach \x/\lab in {0/0,3/s,6/e,12/m} {
    \draw (\x,-0.16) -- (\x,0.16);
    \node[below] at (\x,-0.16) {$\lab$};
  }
  \node[above] at (1.5,0.15) {$\wpre{w}{s}$};
  \node[above] at (4.5,0.15) {delete};
  \node[above] at (9,0.15) {$\wsuf{w}{e}$};
  \draw[<->] (3,1.1) -- (6,1.1)
    node[midway,above] {$q_s=q_e$};
  \node[anchor=east] at (12,1.55) {$\wsuf{w}{e}\pre\wsuf{w}{s}$};
  \draw[->,dashed] (6,-0.8) -- (3,-2.05);
  \draw[->,dashed] (12,-0.8) -- (9,-2.05);
  \node[anchor=east] at (-0.3,-2.3) {$\wpre{w}{m-(e-s)}$};
  \draw[thick] (0,-2.3) -- (9,-2.3);
  \foreach \x in {0,3,9} {\draw (\x,-2.46) -- (\x,-2.14);}
  \node[below] at (1.5,-2.45) {$\wpre{w}{s}$};
  \node[below] at (6,-2.45) {$\wsuf{w}{e}$};
  \node[anchor=west,align=left,font=\footnotesize] at (9.2,-2.3)
    {$q_{m-(e-s)}=q_m$\\$\in\Rej$};
  \draw[dotted] (7,-2.1) -- (7,-3.9);
  \node[above,font=\footnotesize] at (7,-2.05) {$J$};
  \draw[<->] (7,-3.65) -- (9,-3.65);
  \node[below,font=\footnotesize] at (7,-3.85) {first rejected cut};
  \node[anchor=east,font=\footnotesize] at (5.3,-3.8)
    {$J\le m-(e-s)\quad\Longrightarrow\quad e-s\le m-J=r$};
\end{tikzpicture}
\caption{The return gap (\cref{lem:return-gap}). A comparison makes loop deletion
produce the prefix $\wpre{w}{s}\wsuf{w}{e}=\wpre{w}{m-(e-s)}$.
Equal states at $s$ and $e$ preserve the final rejected state, so the shortened
prefix cannot end before $J$. The position of $J$ relative to $s,e$ is schematic;
only $J\le m-(e-s)$ is asserted.}
\Description{A word with cuts s, e, and m, and equal reader states at s and e.
Deleting the shaded interval from s to e moves the suffix to s, yielding a
shorter prefix in the same final rejected state. Its end is at or after the
first rejected cut J, bounding the deleted length by the rejected tail.}
\label{fig:return-gap}
\end{figure}


\subsection{The Depth Bound}

\begin{lemma}[Equal-support width]\label{lem:width}
Let $w$ be a nonempty rejected word with tail $r$, and let $j \le j'$ be levels with
$\Supp_j = \Supp_{j'} \ne \emptyset$. Then $j' - j + 1 \le r+1$.
\end{lemma}

\begin{proof}
For $j=j'$, the bound is immediate. Assume $j < j'$. Since $\Supp_{j'} \neq \emptyset$, all of $\Cores_j, \dots, \Cores_{j'}$ are nonempty.
By \cref{lem:cores}, $t_{i+1} < t_i$ for $j \le i < j'$, so $j' - j \le t_j - t_{j'}$. The state
$q_{t_j}$ lies in $\Supp_j = \Supp_{j'}$, so $q_s = q_{t_j}$ for some $s \in \Cores_{j'}$.
As $\Cores_{j'} \subseteq \Cores_{j+1}$, \cref{lem:cores} shows that $(s,t_j)$ is a comparison,
and \cref{lem:return-gap} gives $t_j - s \le r$. Since $s \le t_{j'}$, we obtain
$j' - j \le t_j - t_{j'} \le t_j - s \le r$.
\end{proof}

\begin{theorem}[Depth from tail]\label{thm:depth}
For every $N$-state reader, each nonempty rejected word with tail $r$ has
$\depth(w) \le 2N(r+1)-1$.
\end{theorem}

\begin{proof}
Let $d_w = \depth(w)$ and $R = r+1$. The supports $\Supp_0 \supseteq \Supp_1 \supseteq
\cdots \supseteq \Supp_{d_w}$ are nested, so the distinct nonempty ones form a strictly
decreasing chain of nonempty subsets of an $N$-element set, and there are at most $N$
of them. By \cref{lem:width}, each of them is the support of at most $R$ levels. Hence at
most $N R$ levels have nonempty support. If every level up to $d_w$ does, then
$d_w + 1 \le N R$. Otherwise, let $j_0$ be the first level with $\Supp_{j_0} = \emptyset$;
then $1 \le j_0 \le N R$, since $\Cores_0$ contains the cuts $0,\dots,\len{w}-1$. All cores
of $u_{j_0}$ are start cores, so $\len{u_{j_0}} \le j_0$ by \cref{lem:cores}. Each further
level up to $d_w$ strictly decreases the length, so $d_w - j_0 \le \len{u_{j_0}} - 1 \le j_0 - 1$,
and $d_w \le 2j_0 - 1 \le 2N R - 1$.
\end{proof}



\Cref{thm:depth} holds for every nonempty rejected word, hole or not; combined with
\cref{thm:brt}, it bounds the depth of some hole by $2N(4N-1)-1$.
